\documentclass[10pt,a4paper]{amsart}
\usepackage[margin=19mm]{geometry}
\usepackage{amsmath,amssymb,mathtools}
\usepackage[scr=rsfs,cal=euler]{mathalfa}
\usepackage{xfrac}
\usepackage[unicode,hidelinks,bookmarks=false]{hyperref}
\newcommand{\ie}{i.e.\ }
\newcommand{\cf}{cf.\ }
\newcommand{\resp}{respectively\ }
\newcommand{\Subsection}{\subsection}
\providecommand{\nobreakdash}{\nobreak\hbox{-}}
\setlength{\emergencystretch}{2em}
\multlinegap0pt
\usepackage{bm}
\usepackage{bbm}
\usepackage{dsfont}
\usepackage{xspace}
\usepackage{cancel}
\usepackage{mathtools}
\usepackage{amsthm}
\makeatletter
\@ifundefined{theorem}{\newtheorem{theorem}{Theorem}}{}
\@ifundefined{lemma}{\newtheorem{lemma}[theorem]{Lemma}}{}
\@ifundefined{observation}{\newtheorem{observation}[theorem]{Observation}}{}
\makeatother
\let\emptyset=\varnothing
\let\setminus=\smallsetminus
\let\backslash=\smallsetminus

\makeatletter
\def\moverlay{\mathpalette\mov@rlay}
\let\setminus=\smallsetminus
\let\backslash=\smallsetminus

\makeatletter
\def\moverlay{\mathpalette\mov@rlay}
\newcommand{\cT}{\mathcal{T}}
\title[Bootstrap percolation of extension hypergraphs]{Bootstrap percolation of extension hypergraphs}
\author{Weichan Liu, Bjarne Sch\"ulke, Xin Zhang}
\date{Source: arXiv:2604.04607v1 (CC BY 4.0)}
\begin{document}
\maketitle

\noindent\emph{Source and licence.} Bootstrap percolation of extension hypergraphs. Version arXiv:2604.04607v1 (CC BY 4.0), distributed under the Creative Commons Attribution 4.0 International licence, as recorded in the arXiv licence metadata for this exact version and shown on the abstract page https://arxiv.org/abs/2604.04607v1. Full rights notice: licenses/arxiv-2604.04607-CC-BY-4.0.txt.\par\smallskip

\noindent\textit{Editorial scope.} Counted unit: the Lemma whose source label is \texttt{lem:good+6} in the article (source0.tex lines 794--796, source label \texttt{lem:good+6}), delivered together with the article's own complete original proof of it (source0.tex lines 798--831, one \texttt{proof} environment of 34 lines). No mathematics is written, repaired or re-derived: statement and proof are the article's own text line for line. Required context retained uncounted: eq:hierarchy (lines 702--704); obs:almostfulldegree (lines 712--717); obs:manycenteredges->good (lines 779--781). Every definition, display and label of those lines is retained unchanged. Omitted from this package and not redistributed: 1--701, 705--711, 718--778, 782--793, 797--797, 832--1226. Nothing inside a retained excerpt is omitted or altered: every retained body is delivered exactly as the article's lines are written, so no omission receipt of any kind accompanies this package. Citations: the retained excerpts carry no citation command at all, so no bibliography entry of the article is transcribed and no reference is redistributed. Reference adaptation: none was needed and none was made; every reference inside the delivered unit resolves inside it, so no unresolved reference or citation is produced. Printed slips kept literal: no printed word, symbol, hypothesis or number of the counted statement or of its proof was altered. Layout and class: the article's own class is \texttt{amsart} with packages including amsmath, amssymb, amsthm, mathrsfs, mathabx, dsfont, bbm, xcolor, hyperref, ocgx2, bookmark, amsrefs, doi, fontenc; the wrapper is the lane's pinned \texttt{amsart} editorial wrapper at 10pt on A4, which restates the theorem-like environments the retained text needs and adds the article's own macro definitions that the retained text needs (\texttt{\textbackslash cT}, \texttt{\textbackslash emptyset}, \texttt{\textbackslash setminus}), with extra wrapper packages (bm, bbm, dsfont, xspace, cancel, mathtools). The delivered numbering is the unit's own. Assets: no retained span contains a figure environment, a graphics-inclusion command, a PDF-page inclusion, a table or a TikZ picture; no image file is copied into this package, the package declares no asset dependency and no image file is redistributed with it. Licence: version arXiv:2604.04607v1 of this article is distributed under the Creative Commons Attribution 4.0 International licence, as recorded in the arXiv licence metadata for this exact version and shown on the abstract page https://arxiv.org/abs/2604.04607v1. A full rights notice is delivered at licenses/arxiv-2604.04607-CC-BY-4.0.txt. Characters: every retained excerpt is pure ASCII, so the delivered engine, XeLaTeX with its default Latin Modern text font, reports no missing character.
	\begin{align}\label{eq:hierarchy}
		t,k \ll \ell \ll c \ll s \ll \Gamma_2 \ll \Gamma_1 \ll \Gamma \ll n_0\,.
	\end{align}

	\begin{observation}\label{obs:almostfulldegree}
		Let~$F$ be a copy of~$F^{(k)}(G)$ in~$H_m$, for some~$m \geq 0$, and let~$uv \in E(C(F))$.
        Further, let~$\mathbf{w} \subseteq V \setminus V(F)$ and~$\mathbf{w}' \subseteq V \setminus V(C(F))$ be~$(k-2)$-sets.
        Then we have~$u\mathbf{w}v\in E(H_{m+1})$ and~$u\mathbf{w}'v\in E(H_{m+2})$.
		%Consequently, for every~$(k-3)$-set~$\mathbf{x} \subseteq V \setminus V(C(F))$, we have~$d_{H_{m+2}}(u\mathbf{x}v)\geq n-t-k+3$.
	\end{observation}

    \begin{observation}\label{obs:manycenteredges->good}
        If~$v\in V$ is~$m$-good for some~$m\in\mathds{N}$, then there is a family~$\cT_v\subseteq V^{(k-2)}$ of pairwise disjoint~$(k-2)$-sets with~$\vert\cT_v\vert=c$ such that~$d_{H_{m+1}}(v\mathbf{u})\geq n-c$ for every~$\mathbf{u}\in\cT_v$.
    \end{observation}

	\begin{lemma} \label{lem:good+6}
		If~$\vert V_m^{\text{good}}\vert\geq s$ for some~$m\in\mathds{N}_0$, then~$\iota\leq m+6$.
	\end{lemma}

	\begin{proof}
		
		We choose pairwise disjoint~$\ell$-subsets~$W_1, \ldots, W_t \subseteq V_m^{\text{good}}$ such that for all integers~$1\leq i<j\leq t$ and for all~$v\in W_i$,~$x\in W_j$, and~$\mathbf{u}\in\cT_v$ we have $v\mathbf{u} x\in E(H_{m+1})$.
		To see that this is possible, first select~$W_1\subseteq V_m^{\text{good}}$ with~$\vert W_1\vert=\ell$ arbitrarily.
		Then assume that~$W_1,\dots, W_{i}$ have been chosen for some~$i\in[t-1]$.
		Note that a vertex~$w\in V_m^\text{good}$ can only not be selected for~$W_{i+1}$ if it lies in~$W_1\cup\dots\cup W_i$ or if there is some vertex~$v\in W_1\cup\dots\cup W_i$ and some~$\mathbf{u}\in \cT_v$ such that~$v\mathbf{u}w\notin E(H_{m+1})$.
        Hence, in light of Observation~\ref{obs:manycenteredges->good}, there are at most~$i\cdot \ell+i\cdot \ell\cdot c^2$ vertices in~$V_m^\text{good}$ which cannot be selected as vertices for~$W_{i+1}$.
		Due to~\eqref{eq:hierarchy}, we can thus still choose~$W_{i+1}\subseteq V_m^{\text{good}}$ as desired.
		
		Arbitrarily pick~$w_i^\star \in W_i$ for each~$i \in [t]$.
		We claim that there are pairwise disjoint~$(k-2)$-sets~$\mathbf{u}_{ij}\in\cT_{w_i^\star}$, for all~$1\leq i<j\leq t$, such that~$\mathbf{u}_{ij}\cap\bigcup_{a\in[t]}W_a=\emptyset$.
		Choose the sets~$\mathbf{u}_{ij}$ greedily following the lexicographic order of~$[t]^{(2)}$, which we denote by~$\preceq_{\text{lex}}$. 
		When picking~$\mathbf{u}_{ij}\in\cT_{w_i^\star}$ at any point during this selection, we only need to guarantee that it is disjoint from the set $$B=\bigcup_{a\in[t]}W_a\cup\bigcup_{i'j'\in[t]^{(2)}\text{ with }i'j'\prec_{\text{lex}} ij}\mathbf{u}_{i'j'}\,.$$
		Since~$\vert B\vert\leq t\cdot\ell+\binom{t}{2}\cdot(k-2)\ll c$ and the sets in~$\cT_{w_i^\star}$ are pairwise disjoint, there is indeed a valid choice for~$\mathbf{u}_{ij}$.
		
		Given~$i,j\in[t]$ with~$i<j$, the choice of~$W_i$ and~$W_j$ ensures that for every~$\mathbf{u} \in \cT_{w_i^\star}$ we have~$w_i^\star \mathbf{u} w_j^\star \in E(H_{m+1})$, whence~$w_i^\star \mathbf{u}_{ij} w_j^\star \in E(H_{m+1})$.
		Consequently,~$H_{m+1}$ contains a copy~$F$ of~$F^{(k)}(K_t^{(2)})$ with~$V(C(F))=\{w_1^\star, \ldots, w_t^\star\}$.
		In fact, since~$w_1^\star,\dots,w_t^\star$ were chosen arbitrarily, for every selection of vertices~$w_i\in W_i$,~$i\in[t]$,~$H_{m+1}$ contains a copy of~$F^{(k)}(K_t^{(2)})$ such that the vertex set of the center of this copy is~$\{w_1, \ldots, w_t\}$.
		By Observation~\ref{obs:almostfulldegree},
		\begin{align}\label{eq:edgescliqueextension}
			\text{$w_i\mathbf{u}w_j$ forms an edge in~$H_{m+3}$}
		\end{align}
		for all~$ij\in[t]^{(2)}$,~$w_i\in W_i$,~$w_j\in W_j$, and every~$(k-2)$-set~$\mathbf{u} \subseteq V \setminus \{w_i,w_j\}$.
		
        Now let~$a,b\in V$ be arbitrary and pick~$w_i\in W_i\setminus\{a,b\}$, for all~$i\in[t-2]$.
		Note that~\eqref{eq:edgescliqueextension} entails that for every~$w_{t-1}\in W_{t-1}\setminus\{a\}$, every~$i\in[t-2]$, and every~$(k-3)$-set~$\mathbf{u'}\subseteq V\setminus\{a,w_i,w_{t-1}\}$ we have~$w_i\mathbf{u'}aw_{t-1}\in E(H_{m+3})$.
		Similarly, for every~$w_t\in W_t\setminus\{b\}$, every~$i\in[t-2]$ and every~$(k-3)$-set~$\mathbf{u'}\subseteq V\setminus\{b,w_i,w_t\}$ we have~$w_i\mathbf{u'}bw_t\in E(H_{m+3})$.
		Using the hierarchy~\eqref{eq:hierarchy}, one can thus observe that it is possible to choose~$\widetilde{\mathbf{u}}_{ij}\subseteq V^{(k-2)}$,~$\mathbf{u}'_i,\mathbf{u}''_i\subseteq V^{(k-3)}$, and~$w_{t-1}^i\in W_{t-1}$ as well as~$w_t^i\in W_t$, for all~$i,j\in[t-2]$, such that the following holds.
        The sets~$\widetilde{\mathbf{u}}_{ij},\mathbf{u}'_i,\mathbf{u}''_i,\{w_{t-1}^i\},\{w_t^i\}$, and~$\{w_1,\dots,w_{t-2},a,b\}$ are all pairwise disjoint and~$w_i\widetilde{\mathbf{u}}_{ij}w_j\in E(H_{m+3})$,~$w_i\mathbf{u}'_iw_{t-1}^ia\in E(H_{m+3})$, and~$w_i\mathbf{u}''_iw_t^ib\in E(H_{m+3})$.
        This means that in~$H_{m+3}$ there is a copy of~$F^{(k)}(K_t^{(2)-})$ whose center has the vertex set~$\{w_1,\dots,w_{t-2},a,b\}$ with~$ab$ being the missing edge.
		Hence,~$H_{m+4}$ must contain a copy of~$F^{(k)}(K_t^{(2)})$ whose center has the same vertex set.
		As argued above (see~\eqref{eq:edgescliqueextension}), this implies that for every~$(k-2)$-subset~$\mathbf{u}\subseteq V\setminus\{a,b\}$ we have~$a\mathbf{u}b\in E(H_{m+6})$.
		Since~$a$ and~$b$ were arbitrary vertices, we infer that~$H_{m+6}$ is complete, whereby~$\iota\leq m+6$.
	\end{proof}

\end{document}
